% =====================================================================================================
% vygotskian_cultural_plane.tex — The Vygotskian machine, part 3: the cultural plane
%                                  (language over the belief hierarchy: culture, sequencing and imagination)
%
% DRAFT (2026-11-19; first drafted as language_hierarchy_section.tex). A proposal section: the architecture it
% describes is NOT yet implemented beyond the parts stated as existing. Source design note:
% docs/LANGUAGE_OVER_BELIEF_HIERARCHY.md.
% Include AFTER vygotskian_machine.tex and vygotskian_internalisation.tex (it refers to their equations and modes).
% Bibliography: vygotskian_machine.bib (shared). Packages: amsmath, amssymb, enumitem, array.
% friston2017deeptemporal duplicates main.tex's friston2017deep under a new key on purpose, so that this file builds
% with vygotskian_machine.bib alone; merge the two keys when the bibliographies are merged.
% Labels keep their original sec:lang_hierarchy:* prefix so that existing cross-references do not break.
% =====================================================================================================

\section{The Cultural Plane: Language over the Belief Hierarchy}
\label{sec:lang_hierarchy}

Section~\ref{sec:inner_speech} treated the language model as an inner voice and asked what language contributes to
thought. In this section we turn to the plane on which, according to Vygotsky, every higher function first appears,
namely the social and cultural plane (Section~\ref{sec:vygotskian}). The vision-language model (VLM) comes to the
robot already containing a large share of human common sense. For the robot, the VLM therefore plays the role of
\emph{culture}: it is an external store of knowledge that extends the robot's mind in the sense of
\cite{clark1998extended}, and the robot inherits this knowledge rather than acquiring it through its own experience
(Section~\ref{sec:inner_speech:why_lm}). The question we address is how such a store can be coupled to a belief state
that is built by probabilistic inference, in such a way that the robot benefits from the knowledge without also
inheriting the fact that the store is neither grounded in the robot's perception nor calibrated.

We propose to connect the VLM to the robot's \emph{belief hierarchy} in three roles, which we call culture,
sequencing and imagination. All three are governed by the same rule, which we state as \textbf{the language model
proposes, the generative models dispose}. The roles are arbitrated by the attention modes of
Section~\ref{sec:inner_speech:modes}, so that the VLM imagines while the automatic level runs well, and monitors and
replans when it does not.

\subsection{The automatic belief hierarchy}
\label{sec:lang_hierarchy:automatic}

The robot builds its belief state without language, using the probabilistic hierarchy of levels 0 to 3 described
in Section~\ref{sec:cortex} (Table~\ref{tab:lh_levels}). One property of that hierarchy is used repeatedly below. A
level-3 metaconcept sends each of its members an empirical prior computed from the arrangement fitted without that
member, which is the cavity distribution of expectation propagation \cite{minka2001expectation},
\begin{equation}
    q_{\setminus i}(\theta) \;\propto\; \frac{q(\theta)}{m_i(\theta)}.
    \label{eq:lh_cavity}
\end{equation}
Here $\theta$ are the parameters of the arrangement (for a kitchen, its axis, worktop height and carcass depth),
$q(\theta)$ is their posterior given all members, and $m_i(\theta)$ is the factor contributed by the posterior of
member $i$. This means that a member is never asked to agree with its own evidence. The hierarchy has three limits,
and the rest of the section addresses them. It cannot say what a room is \emph{for}, because the purpose of a room
is not a geometric fact. It cannot pursue a goal that it was not built to pursue. And it does not look ahead, because
its predictions concern the present scene rather than the outcome of what the robot is about to do.

\subsection{Three roles of the language model}
\label{sec:lang_hierarchy:roles}

\begin{table}[ht]
    \centering
    \small
    \begin{tabular}{|>{\raggedright\arraybackslash}p{0.14\linewidth}|>{\raggedright\arraybackslash}p{0.43\linewidth}|>{\raggedright\arraybackslash}p{0.33\linewidth}|}
        \hline
        \textbf{Role} & \textbf{What the VLM supplies} & \textbf{What decides} \\ \hline
        Culture     & priors and structure that the hierarchy cannot produce itself: room function, class co-occurrence, arrangement schemas, norms & evidence against a null hypothesis (the rules shared by all metaconcepts) \\ \hline
        Sequencing  & goals decomposed into missions over the hierarchy's beliefs, under contextual constraints, with the belief change each step should produce & a mission checker and command gate, the controller, then the belief dynamics \\ \hline
        Imagination & \emph{which} subjunctive questions to put to the robot's own models while the automatic level runs well & the models, run on an imagined copy of the world \\ \hline
    \end{tabular}
    \caption{The three roles. In each, the language model proposes and a generative model disposes.}
    \label{tab:lh_roles}
\end{table}

The three roles rest on the same asymmetry, which was already discussed in Section~\ref{sec:inner_speech:why_lm}.
Probabilistic inference evaluates hypotheses and policies well but cannot generate new ones, whereas a language model
generates them cheaply but cannot evaluate them. Each role therefore combines the two in a generate-and-test loop, in
which the language model generates candidates and the test is carried out by the robot's generative model.

\paragraph{The person is part of the scaffold.}
According to Vygotsky, the knowledge that is internalised comes from a more knowledgeable other, and the support this
other provides, known as scaffolding \cite{wood1976role}, is withdrawn as the learner's competence grows. The VLM is
one such other, since it makes culture available on demand. The person living with the robot is the second one, and
already plays this role in the deployed system. The person sets goals, answers the questions that the robot cannot
settle by looking, confirms a claim before it is stored, and approves a mission before it starts. In the cultural
plane, these answers are the most valuable evidence the robot receives. A room function confirmed by the person
contributes to the counts $n_f$ in \eqref{eq:lh_beta_prior}, and a plan that the person approved and that succeeded
contributes to the habit term $E$ in \eqref{eq:lh_habit}. The decision to ask the person rather than to look is
itself assigned a cost, as the language agent already does for objects that may have moved. Both forms of scaffolding
are designed to be withdrawn. The weight of the VLM decreases as the pseudo-count $\kappa\beta_f$ is outweighed by
experience, and the person is asked less often as the robot's own counts grow.

\subsection{Culture: a level-4 function hypothesis}
\label{sec:lang_hierarchy:culture}

\paragraph{Latent and evidence.}
For each room $r$ and function $f$ we maintain an independent hypothesis ``$r$ serves $f$'', which is compared with
the null hypothesis ``$r$ is a generic space''. The hypotheses are independent rather than mutually exclusive
categories, because open-plan rooms serve several functions, and because this is how the shared metaconcept
machinery already treats its instances. The evidence for a hypothesis comes from one observation model. For each
class $c$, the observation is a binary variable $z_c$ that equals one if $c$ has been detected in $r$, so that ten
chairs do not count as ten separate pieces of evidence. Under $f$, the class is present with probability
$\pi_{cf}$ and, if present, is detected with probability $d_c$, so that $P(z_c = 1 \mid f) = \pi_{cf} d_c$. Under
the null hypothesis the same holds with $\pi_{c0}$ in place of $\pi_{cf}$. The log-likelihood ratio of one class is
therefore $\log(\pi_{cf}/\pi_{c0})$ when it has been detected, and $\log\big((1 - \pi_{cf} d_c)/(1 - \pi_{c0}
d_c)\big)$ when it has not, and the detectability cancels from the first case. A detection may itself be false, so
we take the expectation of the log-ratio under the probability $p_c$ that the class has truly been detected,
\begin{equation}
    \Lambda_f(r) \;=\; \log\frac{P(f)}{P(0)} \;+\; \sum_c \Big[\, p_c \log\frac{\pi_{cf}}{\pi_{c0}}
    \;+\; (1 - p_c)\log\frac{1 - \pi_{cf}\, d_c}{1 - \pi_{c0}\, d_c} \Big],
    \qquad p_c = 1 - \prod_{i \in c} (1 - p_i).
    \label{eq:lh_function_evidence}
\end{equation}
Here $p_i$ is the existence probability of instance $i$ of class $c$ in $r$, the product runs over the instances
that have been observed at least once, and $p_c = 0$ for a class with no instance. The first term is the prior
log-odds of the function, which is elicited from the VLM together with $\pi_{cf}$. The probabilities $\pi_{c0}$ of
the null hypothesis are elicited in the same way, by asking the VLM about a generic room, and they are replaced by
the robot's counts over all the rooms it has confirmed. Put simply, a class that was found speaks for the functions
that expect it, and a class that was looked for and not found speaks against them in proportion to how thoroughly it
was looked for. When the presence of every class is certain, the sum is the exact log-likelihood ratio. When it is
not, the expectation of the log-ratio is an approximation to the log of the marginal likelihood, which we accept
because it keeps the contribution of each class separate.

Two dependences must be handled. First, the existence probability $p_i$ of an instance already depends on whether the
instance was seen again where it should have been, but it does not depend on the search for other instances of the
class. The detectability $d_c$ therefore enters the evidence once, through the second term. Second, several classes
are not independent given the function. A counter, a countertop, a hood and the cabinets of a kitchen occur together.
Classes that are members of a level-3 aggregate therefore contribute through the aggregate, as a single observation,
and not one by one. The dependence that remains between classes outside any aggregate makes $\Lambda_f$ somewhat
overconfident, and we list it among the limitations. At present $d_c$ is the room's exploration completeness
multiplied by a per-class detection rate. This scalar cannot register a search directed at a particular class, and
it must be replaced by a detectability model computed from the viewpoints the robot has actually occupied before the
second stage described below is enabled. The quantity $\Lambda_f(r)$ is the existence log-odds of the hypothesis,
and the birth and death of hypotheses follow the rules of the metaconcept machinery.

\paragraph{Priors from culture, replaced by experience.}
The conditional probabilities $\pi_{cf}$ cannot be supplied by any geometric model. We elicit them from the VLM over
a closed vocabulary, consisting of the class names already used in the graph and a fixed set of functions, and give
each of them a Beta prior whose strength must be earned through measured calibration,
\begin{equation}
    \pi_{cf} \sim \mathrm{Beta}\big(\alpha + \kappa\beta_f\,\pi^{\text{VLM}}_{cf},\;
    \alpha + \kappa\beta_f\,(1-\pi^{\text{VLM}}_{cf})\big),
    \qquad
    \hat\pi_{cf} = \frac{\alpha + \kappa\beta_f\,\pi^{\text{VLM}}_{cf} + n_{cf}}{2\alpha + \kappa\beta_f + n_f}.
    \label{eq:lh_beta_prior}
\end{equation}
Here $\alpha > 0$ is a small base count that keeps the prior defined when $\beta_f = 0$ and $n_f = 0$, and $\kappa$
is a small pseudo-count, chosen so that the VLM carries the weight of a few observed rooms rather than a hundred.
The factor $\beta_f \in [0,1]$ is the robot's estimate of the VLM's reliability for $f$, defined as in
\eqref{eq:inner_prior}. It is the posterior mean of a $\mathrm{Beta}(1,1)$ distribution over the probability that,
in a room whose function a person has stated, the hypothesis seeded by the VLM agrees with the person, and it is
updated only on such rooms, so that the VLM is never calibrated against labels that its own priors produced. The
count $n_f$ is the number of distinct rooms confirmed as serving $f$, and $n_{cf}$ is the number of those rooms that
contained $c$. A room counts once, however often it is visited. This is the same Dirichlet form as
\eqref{eq:inner_prior}, restricted to two outcomes. The effective weight of the VLM, $\kappa\beta_f/(\kappa\beta_f +
n_f)$, is large while few rooms have been confirmed and decreases as experience accumulates.

The counts must not be allowed to confirm the prior that produced them. If a room were counted because a posterior
computed with $\hat\pi_{cf}$ had become confident, and the classes it contained then increased $n_{cf}$, the VLM's
prior would confirm itself through the counts, without any new information entering. We therefore admit a room into
$n_f$ in one of two ways only. Either a person has stated its function, or the evidence \eqref{eq:lh_function_evidence}
for $f$ remains positive when it is recomputed with the VLM's pseudo-counts removed and with the room's own
contribution removed from the counts, that is, with a prior built only from the other rooms the robot has
confirmed. The second route is closed until a few rooms have been confirmed by a person. The same elicitation will
replace priors that are currently written by hand. For example, the class-membership
log-odds of the kitchen metaconcept are documented in its configuration as ``elicited, not measured'' and are still
awaiting a corpus from which to measure them.

\paragraph{How a function hypothesis influences lower levels without confirming itself.}
In the first stage, a function hypothesis sends nothing into the beliefs of its members. In the second stage it sends
only exploration targets, namely classes that are expected but not yet seen, each with a coarse placement. The
exploration drive and the next-best-view planners use these targets to value candidate viewpoints, and never treat
them as evidence that an object exists. This introduces a selection bias in where the robot looks, but it does not
contaminate the evidence, provided that the detectability in \eqref{eq:lh_function_evidence} is computed from the
viewpoints that the robot actually occupied, because absence is then counted only in the places that were searched. Only at a later stage, and only if its record justifies it, would a function seed
level-3 births, and in that case the down-date \eqref{eq:lh_cavity} would apply to the members it seeded, so that
an object born because ``kitchen'' predicted it is not counted as evidence for ``kitchen''. The cavity rule covers
only this route. It does not prevent self-confirmation through the counts of \eqref{eq:lh_beta_prior}, which is
prevented instead by the admission rule stated with that equation, and it does not protect against a VLM that is
wrong in the same way about every room, which only rooms confirmed by a person can reveal.

\paragraph{Norms.}
Each function also carries the contextual norms that the VLM supplies for it, such as moving quietly and slowly in a
bedroom at night, keeping clear of the work area in a kitchen that is in use, and never blocking a doorway. These
norms are not evidence. They are preferences of the generative model, and
Section~\ref{sec:lang_hierarchy:sequencing} uses them as planning constraints.

\paragraph{Model-level surprise.}
Levels 2--3 already produce signals indicating that their \emph{model} is wrong, as opposed to merely uncertain.
Such signals include an arrangement whose Bayes factor collapses, a member with a large cavity residual, a class that
persistently receives a membership score near zero, and two competing schemas that are both weak. These signals
trigger diagnosis in focused mode. An example from the live system illustrates this. A refrigerator received a
kitchen membership of $0.00$ because only base and wall units were modelled, so that its top, at $1.85$\,m, was
compared with a worktop plane at $0.90$\,m. A person diagnosed the structural defect, and adding a third, tall tier
to the model corrected it. In focused mode, the task of the VLM is to propose revisions of this kind, which change
the model itself rather than its parameters. For now, accepting such revisions remains the task of a person.

\subsection{Sequencing over the hierarchy}
\label{sec:lang_hierarchy:sequencing}

\paragraph{Goals at their own level.}
In the deployed system, mission predicates already refer to beliefs, for example a room explored with probability
$\geq 0.9$, a door open with probability $< 0.5$, or an object class that has been seen. The hierarchy extends the
vocabulary in which goals are stated and decomposed. This corresponds to the deep temporal structure of Active
Inference \cite{friston2017deeptemporal}, in which a goal is planned at its own level and then refined downward
(Table~\ref{tab:lh_goals}).

\begin{table}[ht]
    \centering
    \small
    \begin{tabular}{|>{\raggedright\arraybackslash}p{0.2\linewidth}|>{\raggedright\arraybackslash}p{0.27\linewidth}|>{\raggedright\arraybackslash}p{0.43\linewidth}|}
        \hline
        \textbf{Goal stated at} & \textbf{Example} & \textbf{Decomposition} \\ \hline
        Level 4, function    & ``go where people eat''           & room with the highest $\Lambda_{\text{dining}}$ $\to$ its dining set $\to$ a standpoint \\ \hline
        Level 3, aggregate   & ``check every chair at the table'' & the ring's members \emph{and its empty slots}, predicted but not yet seen \\ \hline
        Level 2, object      & ``approach the fridge''           & an approach to the object \\ \hline
        Level 1, room        & ``go to room 2''                   & the doors to cross, from the spatial memory of the rooms already visited \\ \hline
        Level 0, self        & ``turn 90 degrees''               & a relative motion \\ \hline
    \end{tabular}
    \caption{Goals are planned at their own level of the hierarchy and refined downward.}
    \label{tab:lh_goals}
\end{table}

\paragraph{Epistemic sequences.}
The exploration drive is the robot's compiled epistemic plan, that is, a fixed routine for gathering information, and
it is the right default for building the map. When the robot has a specific goal, however, the drive becomes
inefficient. To find a mug in a partly known house, the VLM writes a sequence that the drive cannot produce. The
robot goes to the room with the highest kitchen evidence (level 4), looks along the worktop plane where mugs are
usually placed (level 3), and approaches a mug once one is seen. The hierarchy indicates where the information is,
and the VLM determines the order in which to gather it and when to stop.

\paragraph{Every step predicts its belief outcome.}
At present, a mission step reports only completion or failure. We require each step also to declare the change in
belief that it should produce. A door crossing predicts that the robot's current room changes and that the new room
is recognised. An exploration step predicts that completeness rises and that some $\Lambda_f$ moves away from zero.
An approach predicts that the existence probability of the target stays high as the distance to it decreases. Let
$\phi_k$ be the predicate over beliefs that step $k$ predicts, which must hold by the end of the step, and let
$\psi_j$ be the assumptions on which the running plan relies (``door 3 is open'', ``room 2 is still the
kitchen''). Both enter the demand signal \eqref{eq:inner_demand} as surprisal terms,
\begin{equation}
    s^{\text{plan}}_t \;=\; -\log P_t(\phi_k) \;\;\text{at the step's deadline},
    \qquad
    s^{\text{assume}}_{j,t} \;=\; -\log P_t(\psi_j) \;\;\text{continuously},
    \label{eq:lh_plan_surprise}
\end{equation}
where $P_t$ is the current belief. The first term detects a step whose outcome diverged from its prediction even
though the step reported success. The second term raises attention when an assumption weakens, before any step has
failed, in the same way as a driver slows down on seeing a school sign rather than braking only when a child appears.
In either case, the focused voice replans from the belief state that the robot actually has. A plan thereby becomes
the same kind of object as an inner question (Section~\ref{sec:inner_speech:questions}), namely a proposal whose
predictions are confirmed or refuted by the belief dynamics.

\paragraph{Constraints from context.}
The norms attached to the active function hypotheses become invariants that apply to the whole mission, such as
avoiding a region, limiting the speed or keeping a distance. They are monitored for violation throughout the mission
rather than checked only once at its start.

\paragraph{Habits.}
Sequences that recur and succeed are compiled into habits. In Active Inference, the prior over policies combines a
habit term $E$ with the expected free energy $G$, weighted by a precision $\gamma$
\cite{friston2017process, parr2022active},
\begin{equation}
    p(\pi \mid g, x) \;=\; \sigma\big(\ln E(\pi \mid g, x) - \gamma\, G(\pi)\big),
    \qquad
    G(\pi) \;=\; -\sum_k \mathbb{E}\big[\text{information gain of } \phi_k\big]
    \;-\; \mathbb{E}\big[\ln P(\text{goal } g \text{ holds})\big].
    \label{eq:lh_habit}
\end{equation}
Here $g$ is the goal type, $x$ the context given by the function and arrangement hypotheses of the current room, and
$\sigma$ the softmax over the candidate missions. The habit term $E$ is a Dirichlet count of the missions that
succeeded for that goal type in that context. The expected free energy of a mission is computed from its own declared
step outcomes $\phi_k$, as the information that the steps are expected to gain plus the expected log-preference for
the goal, and it is evaluated with the imagination tools of Section~\ref{sec:lang_hierarchy:imagination}. Put
simply, a mission is preferred when it has worked before in situations of the same kind and when it is expected to
reach the goal while learning what the robot needs to know.

The VLM does not enter $G$ directly. It defines the set of candidate missions over which the softmax is taken, and
the norms it supplies enter the preferences. A habit can run without the VLM because of the way it is stored. When a
mission succeeds, its references to particular objects and rooms are replaced by the roles that the hierarchy
assigned to them, so that ``go to chair 34 beside table 2'' is stored as ``go to a chair of the dining set in the
current room''. The stored mission is then retrieved by goal type and context, and its roles are bound to the
objects of the new situation by the arrangement agents, without any call to the language model. The VLM is still
needed to understand what a person asks for, but not to plan a mission of a familiar type. In this sense compiled
habits form plan schemas defined over arrangement schemas. Each habit also carries the assumptions $\psi_j$ under
which it succeeded, so that the robot can detect when the habit no longer applies. Storing a validated mission and
retrieving it is a form of caching, and we do not claim more for it. What distinguishes it from a plain cache is that
retrieval is weighed against expected free energy and that the stored assumptions are monitored.

\subsection{Imagination while the automatic level runs well}
\label{sec:lang_hierarchy:imagination}

\paragraph{Principle.}
In this framework, imagining means running the robot's own generative models offline on an imagined copy of the
world. Both of these components already exist in the deployed language agent, as tools that a person can query. They
comprise a copy of the room's facts, which interventions can edit but which cannot write to the live graph; the
controller's own model of whether the robot fits through a gap (``would I fit through door 1 if it were half
open?''); the perception models (``would I see the fridge from there?''); concept priors that are sampled and
rendered through the camera and LiDAR models with a Monte-Carlo estimate of the probability of detection (``how do
you imagine a chair?''); and the odds of success for a route, estimated from the crossings recorded that day. In
wandering mode the robot puts these questions to itself. The VLM chooses what to imagine, that is, the seed, the
intervention and the question, and the models compute the answer. If the VLM answered a subjunctive question by
itself, it would be confabulating rather than imagining.

\paragraph{Five kinds of imagining.}
Table~\ref{tab:lh_imagination} lists the five kinds according to what they prepare. The first two serve the near
future, the next two the recent past, and the last one the model itself.

\begin{table}[ht]
    \centering
    \small
    \begin{tabular}{|>{\raggedright\arraybackslash}p{0.16\linewidth}|>{\raggedright\arraybackslash}p{0.21\linewidth}|>{\raggedright\arraybackslash}p{0.27\linewidth}|>{\raggedright\arraybackslash}p{0.25\linewidth}|}
        \hline
        \textbf{Kind} & \textbf{Seed} & \textbf{Question to the models} & \textbf{Product} \\ \hline
        Rehearsal & next steps of the running mission & success odds per leg; visibility from the next standpoint; fit through a partly open door & a contingency cache keyed by the assumptions that would trigger each alternative \\ \hline
        Scene completion & empty slots; expected but unseen classes & sample the class at its predicted place; probability of detection from candidate viewpoints & viewpoints with their detection probability \\ \hline
        Counterfactual replay & a failed or slow episode & the same episode through another door, or with another opening & lessons as testable hypotheses; assumptions for habits \\ \hline
        Anomaly pre-analysis & an unexplained residual; a question closed as ``unknown'' & candidate classes rendered at that place against the detection record & ranked hypotheses for the next focused episode \\ \hline
        Structural replay & logged metaconcept inputs & refit under an alternative schema or tier & structural proposals for review \\ \hline
    \end{tabular}
    \caption{Imagination in wandering mode. The VLM chooses the seed and the question; the robot's models answer.}
    \label{tab:lh_imagination}
\end{table}

Structural replay is the kind of process that could have found the tall tier, by refitting the logged kitchen
memberships with one additional tier and observing that the evidence rises.

\paragraph{What to imagine first.}
Imagination competes for a bounded computational budget and can be pre-empted at any moment, so seeds are ranked by
the value of their product for upcoming decisions. Following the account of prioritised replay in
\cite{mattar2018prioritized}, the value of a seed $s$ is
\begin{equation}
    V(s) \;=\; \mathrm{need}(s)\,\cdot\,\mathrm{gain}(s),
    \label{eq:lh_priority}
\end{equation}
where the need measures how soon the situation that $s$ prepares for will occur, and the gain measures how much the
product could change the decision taken in that situation. Need is highest for the next legs of the running mission,
followed by the rooms the drive will enter next, then open questions, and finally old failures. Gain is highest for
legs with uncertain odds of success, expectations whose detection probability is near one half, and hypotheses that
are currently tied.

\paragraph{Where the results of imagination go, and how much they are trusted.}
A contingency cache is attached to the running mission, so that when the robot returns abruptly to focused mode, the
focused voice verifies an alternative that has already been prepared instead of composing a new one. Viewpoints are
passed to the exploration drive with bounded weight, through the same channel as the targets of the function
hypotheses. Ranked hypotheses become priors for the next focused episode. Lessons become candidate habits, which are
trusted only after later executions confirm them. Every imagined episode is stored with a provenance that marks it as
imagined, and is never written to the live graph or to factual episodic memory
(Section~\ref{sec:inner_speech:modes}). Every product that makes a prediction is scored once the corresponding
observations arrive, so that each kind of imagination earns trust in exactly the same way as the priors of the VLM
do, and a kind whose predictions keep failing loses weight.

\paragraph{Gating.}
Imagination runs only in wandering mode. It is cancelled as soon as the demand signal crosses $\theta_{\text{hi}}$,
and it resumes if its seed is still relevant once demand has subsided. Extended reasoning by the VLM is permitted in
this mode and in no other (Section~\ref{sec:inner_speech:cot}).

\subsection{The three roles in the three modes}
\label{sec:lang_hierarchy:modes}

\begin{table}[ht]
    \centering
    \small
    \begin{tabular}{|>{\raggedright\arraybackslash}p{0.13\linewidth}|>{\raggedright\arraybackslash}p{0.23\linewidth}|>{\raggedright\arraybackslash}p{0.16\linewidth}|>{\raggedright\arraybackslash}p{0.2\linewidth}|>{\raggedright\arraybackslash}p{0.16\linewidth}|}
        \hline
        \textbf{Mode} & \textbf{Entered when} & \textbf{Culture} & \textbf{Sequencing} & \textbf{Imagination} \\ \hline
        Focused    & model-level surprise; step divergence or weakened assumption \eqref{eq:lh_plan_surprise}; the person speaks & diagnose the model; propose a revision & replan from the actual belief state, contingency cache first & --- \\ \hline
        Monitoring & between thresholds & --- & expected outcomes and assumptions checked; no VLM calls & --- \\ \hline
        Wandering  & demand low for a dwell time & offline elicitation; schema proposals & drafts for goals ahead & the five kinds, ranked by \eqref{eq:lh_priority} \\ \hline
    \end{tabular}
    \caption{Arbitration of the three roles by the attention modes of Section~\ref{sec:inner_speech:modes}.}
    \label{tab:lh_modes}
\end{table}

The cultural plane therefore adds two kinds of input to those of the demand signal listed in
Section~\ref{sec:inner_speech:modes}, namely model-level surprise from levels 2--3 and the plan surprisals
\eqref{eq:lh_plan_surprise}. The three roles form a closed loop.
Imagination prepares what focused mode will need, focused mode produces the episodes that imagination later replays,
and both supply the counts with which the robot's experience replaces the culture of the VLM in
\eqref{eq:lh_beta_prior} and with which the robot's habits replace the plans of the VLM in \eqref{eq:lh_habit}. The
share of the robot's competence that depends on the VLM should therefore shrink with experience for routine matters,
and remain for situations that are novel, social or structural.

\subsection{Relation to prior work}
\label{sec:lang_hierarchy:related}

\paragraph{Language models as robot planners.}
Language models have been used to sequence robot skills, with each candidate skill scored by a learned estimate of
whether it can succeed in the current state \cite{ahn2022saycan}; to re-plan from textual feedback about success, the
scene and human input \cite{huang2022inner}; to plan over large three-dimensional scene graphs by searching a collapsed
version of the graph and verifying the plan in a simulator \cite{rana2023sayplan}; and to bias object search with soft
common-sense constraints \cite{zhou2023esc}. Other work has asked the language model to write the robot's program
directly, either as code that calls perception and control functions \cite{liang2023code} or as a situated program
whose prompt lists the actions and objects available in the current environment \cite{singh2023progprompt}. In all of
these systems the language model produces the plan, and the robot's role is to execute it and to report success or
failure. In the Vygotskian machine the language model also writes plans, but each step must state the change in belief
it is expected to produce, and the plan is checked against the belief state while it runs.

\paragraph{Scene graphs.}
Three-dimensional scene graphs organise a building into layers for the building, its rooms, its objects and the
cameras that observed them \cite{armeni2019scene}. Later work added a layer of places and built such graphs in real
time from a moving robot \cite{hughes2022hydra}, and with
open-vocabulary labels obtained from vision-language models \cite{gu2024conceptgraphs}. These graphs are the closest
relative of the belief hierarchy of Section~\ref{sec:cortex}. The difference lies in what a node is. In these systems
a node is an estimate that a mapping pipeline has committed to, whereas in \cortex{} every node is a posterior, with a
covariance and an existence probability, maintained by an agent that keeps testing it against new evidence. This is
what allows a level to send priors back down to the level below it, and what allows the cultural level to be
introduced as one more hypothesis that has to earn its place by evidence.

\paragraph{Semantic mapping and room categories.}
Inferring the category of a room from the objects it contains is an established problem in semantic mapping. Pronobis
and Jensfelt combined object observations, room appearance and default knowledge about which objects occur in which
kinds of room in a single probabilistic model \cite{pronobis2012large}. The function hypothesis of
Section~\ref{sec:lang_hierarchy:culture} belongs to this line of work, and it differs in three respects: its default
knowledge is elicited from a vision-language model rather than collected by hand, its prior is gradually replaced by
counts from rooms whose function has been confirmed, and an expected object that has not been found counts as evidence
only in proportion to how thoroughly the robot has searched for it.

\paragraph{Inner speech in robots.}
Mirolli and Parisi \cite{mirolli2011towards} and Colas et al.\ \cite{colas2022language} have argued for robots and
artificial agents that use language as a tool for thought (Section~\ref{sec:vygotskian}). Pipitone and Chella built a
robot that speaks its reasoning aloud while it cooperates with a person, and found that hearing this inner voice made
the robot's decisions more transparent to its partner \cite{pipitone2021robots}. In their work inner speech is a
running commentary on deliberation. In the Vygotskian machine it is an epistemic action, recruited only when the
automatic level is surprised, and its conclusions are accepted only when perception confirms them.

\paragraph{Active Inference, language and amortisation.}
Active Inference has been applied to robot perception, estimation and control \cite{lanillos2021active}. Pezzulo et
al.\ have used it to distinguish generative models that are acquired by acting in the world from those that are
trained passively on its records, and have argued that large language models are of the second kind and therefore lack
the grounding that action provides \cite{pezzulo2024generating}. The Vygotskian machine takes this distinction as a
design principle: it couples a passively trained model to an agent that acts, and lets only the acting agent decide
what is believed. Developmental robotics has long studied how a robot can acquire competence in stages, as children do
\cite{cangelosi2015developmental}, and cognitive science describes how a learner can replace expensive inference with
fast learned mappings that approximate it, a process known as amortised inference \cite{gershman2014amortized}.
Internalisation, as defined here, joins the two. Knowledge first borrowed from the language model is used through
explicit inference and language, and only what has been confirmed by the robot's own experience is compiled into
priors and habits that work without language.

\paragraph{Summary of the differences.}
The present proposal differs from this body of work in three respects. First, the language model is coupled to a
\emph{probabilistic} hierarchy rather than to a symbolic or deterministic representation, and its contributions enter
this hierarchy as priors with measured calibration that are admitted only through evidence against a null hypothesis.
Second, its plans are predictions over beliefs, and a divergence from these predictions acts as a surprise signal,
whereas in earlier work plans are scripts that are checked only for success. Third, its three roles are arbitrated by
an attention mechanism that leaves the language model free to imagine when the automatic level is sufficient.

\subsection{Testable predictions}
\label{sec:lang_hierarchy:predictions}

Each prediction is tested by comparing one policy against another inside the same closed loop, with the body, the
sensors and the environment held fixed. Section~\ref{sec:experiments} assigns each prediction to an experiment.
\begin{enumerate}[label=\textbf{L\arabic*}, start=0]
    \item VLM-elicited class priors are rank-correlated with the hand-authored kitchen membership priors and are
    better calibrated than a prior that always answers ``likely'' against class frequencies measured in a labelled
    corpus of real homes, and substituting them changes no correct live membership decision.
    \item A read-only function hypothesis recovers room functions on a labelled corpus of real homes better than
    uniform, hand-written and co-occurrence priors; in the simulated apartment, each $\Lambda_f$ leaves its prior
    log-odds only as exploration proceeds; removing the refrigerator weakens the kitchen hypothesis without
    destroying it; and a deliberately wrong elicited prior is overridden by the robot's counts.
    \item With predicates and goals at levels 3--4, goal-directed search (``find a mug'') takes less time and distance
    than the exploration drive alone, and fewer objects remain unseen in rooms counted as explored.
    \item With per-step expected outcomes and assumptions in the demand signal, a larger fraction of the seeded
    divergences (a door shut mid-mission, an object moved) is caught before any step fails, and it is caught earlier,
    than when the demand signal sees only step completion and failure.
    \item Rehearsal and scene completion reduce the latency from surprise to a validated new plan, and the number of
    views needed to confirm an expected object.
    \item Counterfactual and structural replay propose the missing tall tier when the two-tier kitchen model is
    re-created.
    \item As confirmed rooms and successful missions accumulate, the estimate $\beta_f$ converges, the effective weight
    of the VLM in \eqref{eq:lh_beta_prior} falls, seeded and unseeded posteriors agree
    on held-out evidence, and VLM calls per hour fall while task success holds.
\end{enumerate}
If L1 fails, the cultural layer is abandoned before any coupling is attempted. If L4 fails, imagination does not
justify its computational cost, and the robot should simply remain idle when the automatic level is sufficient.

\subsection{Limitations}
\label{sec:lang_hierarchy:limitations}

The function evidence \eqref{eq:lh_function_evidence} takes the expectation of the log-likelihood ratio rather than
the log of the marginal likelihood, and it treats classes outside any aggregate as independent given the function,
which makes it somewhat overconfident. The pseudo-count $\kappa$ and the reliability estimate $\beta_f$ trade early usefulness against overconfidence.
Moreover, the VLM's notion of a typical home carries a cultural bias that only the robot's own counts can correct,
although the log-odds formulation down-weights unexpected classes without ever excluding them. Per-step expected
outcomes require each mission operation to know its own postcondition over beliefs, which is straightforward for door
crossing and approach but less so for open-ended exploration. A single VLM serves the person, focused replanning and
imagination, so strict pre-emption limits how much a busy robot can imagine. Finally, structural revisions are
proposed by the VLM but applied by a person, and allowing the robot to modify its own generative models is outside
the scope of this proposal.
